Nguồn: \href{http://www.vanhoanghean.com.vn/component/k2/31-nguoi-xu-nghe/12442-giao-su-phan-dinh-dieu-tai-nang-va-tam-huyet}{Báo Văn Hoá Nghệ An}

Tôi biết GS. Phan Đình Diệu từ ngày còn học trường cấp III Phan Đình Phùng, Hà Tĩnh, năm học 1953-54. Hồi đó học sinh phải học ban đêm để tránh máy bay Pháp bắn phá, tôi lại học lớp khác, nên chưa biết anh Diệu, chỉ nghe đồn anh rất giỏi toán. Thế rồi một lần trên đường từ trường, ở xã Ngu Lâm, huyện Đức Thọ về nhà, tôi gặp anh Diệu chống gậy vừa đi vừa dò đường bị ngập nước sau trận lụt. Chúng tôi trở thành bạn đường, vừa đi vừa trò chuyện, cho tới nhà anh ở Can Lộc thì chia tay, tôi đi tiếp về Thạch Hà. Đến giữa năm học này, các bên ký hiệp định Geneve, hoà bình được lập lại. Lúc này, Ty giáo dục chủ trương chuyển một số học sinh các lớp 9 trường cấp III Phan Đình Phùng ở Ngu Lâm về thị xã Hà Tĩnh làm thành như một phân hiệu. Vì thế tôi được học cùng lớp với anh Diệu. Năm 1955, tình cờ tôi gặp lại anh trên một chuyến xe khách từ Vinh ra Hà Nội để thi đại học. Phải chuyển xe hai ba lần xe, đường lại khó đi, vì thế chúng tôi chẳng còn tâm can đâu để trò chuyện với nhau. Từ đó, anh ở lại Hà Nội, học trường Đại học sư phạm, tốt nghiệp xuất sắc, được giữ làm cán bộ giảng dạy, rồi sang Liên Xô bảo vệ luận án phó tiến sĩ và tiến sĩ toán học. Tôi vì dạy học ở xa, nên chỉ hỏi thăm và được bạn bè cũ cho biết vài thông tin về anh.

Khoảng năm 1970, ngành giáo dục Thanh Hóa tổ chức một Hội nghị khoa học giáo dục và tôi được giao nhiệm vụ rà soát, chọn lọc các bản báo cáo, phần lớn là sáng kiến, kinh nghiệm, để trình bày trong Hội nghị. Đặc biệt có báo cáo “Cách giải bài toán Fermat” của giáo viên tổ tự nhiên một trường cấp 3 lớn nhất tỉnh. Được lãnh đạo Ty cử ra Hà Nội xin ý kiến nhận xét, đánh giá của các nhà khoa học, tôi nghĩ ngay đến Tiến sĩ toán học Phan Đình Diệu, bạn học phổ thông cũ, nghe nói đang công tác ở Uỷ ban KHKT Nhà nước và liền vội tới nơi làm việc tìm gặp. Nhưng anh đang bận thuyết trình một vấn đề về tin học tại hội trường Uỷ ban. Tôi ngồi vào chỗ phía cuối, vì không hiểu tin học là gì, nên chủ tâm không phải ngồi nghe, mà chỉ nhằm mục đích chờ gặp anh. Trong hội trường có mặt cả mấy nhà khoa học nổi tiếng như GS. Tạ Quang Bửu, VS. Trần Đại Nghĩa. Khoảng hơn 11 giờ, buổi thuyết trình kết thúc. Tuy xa nhau hàng chục năm, song tôi nhận ra anh Diệu ngay; người thấp, mập, trán cao và dọng nói không thể nào nhầm với người khác. Anh ra sau cùng, song mấy thính giả trẻ tuổi vẫn đi theo, như thể tranh thủ trao đổi hay hỏi thêm vấn đề gì đó. Nhìn thấy tôi, anh có vẻ ngạc nhiên. Anh vốn có cá tính từ thuở thiếu thời là ít bộc lộ cảm xúc ra ngoài, nên người không hiểu khi gặp thường nhầm tưởng anh có thái độ lạnh lùng, dửng dưng. Đoán trong đầu anh đang đặt câu hỏi sao bao năm xa cách nay bổng nhiên gặp lại, nên tôi đi ngay vào vấn đề. Anh rủ tôi cùng về nhà chơi rồi hẳng hay. Hồi này, gia đình anh đang ở khu tập thể Giảng Võ. Tuy đã trưa, cơm nước dọn sẵn, cả nhà đang chờ, nhưng anh vẫn bảo tôi đưa xem bản báo cáo. Đọc xong, anh nhận xét ngay, định lí Fermat đã được giải rồi, nhưng phương pháp đơn giản hơn nhiều. Anh nói thêm: ở địa phương thiếu sách vở tham khảo, nên cũng cần khen thưởng cho các tác giả. Thầm nghĩ, không biết ông bạn mới đọc qua có gì nhầm không vì ở nhà đánh giá báo cáo là một “ phát minh”, tôi đề nghị anh cho xem cuốn sách có lời giải bài toán. Anh rút ngay trên giá một cuốn sách chữ Nga, chỉ cho tôi xem và giải thích sự đơn giản của cách giải trong sách. Nể vì đã bắt cả gia đình bạn phải chờ, nên khi nghe cô Hương vợ anh thiết tha mời cùng ăn, tôi không thể từ chối, tuy biết còn thiéu suất cơm của tôi. Bữa ăn của một nhà khoa học hồi này cơm độn khoai khô và thức ăn chỉ có một đĩa rau muống luộc, bát nước canh rau và mấy con cá rô kho mặn.

 Tôi được biết cũng khoảng thời gian trên, sau lần đi Pháp về, anh Diệu đã phát hiện triển vọng to lớn của một lĩnh vực khoa học rát mới mẽ là tin học, khoa học máy tính và anh đã miệt mài nghiên cứu. Năm 1977, theo đề nghị của anh, chính phủ cho thành lập Viện khoa học tính toán và điều khiển học. Anh được cử làm Viện trưởng đồng thời giữ chức Viện phó Viện khoa học Việt Nam. Năm 1980, anh được phong hàm Giáo sư. Một lần tôi cùng anh về quê Hà Tĩnh, có cả cháu Hà khoảng 10 tuổi, con gái thứ của anh. Trên đường chuyện trò, anh cho biết ở một số nước phương Tây người ta tuân thủ luật giao thông hết sức nghiêm túc. Có lần chiếc xe anh đi dừng ở một địa điểm quá qui định ít phút, cảnh sát chẳng cần gặp lái xe kiếm tra giấy tờ, mà dán vào xe một phiếu phạt. Lái xe nhìn thấy, lập tức đi nộp tiền phạt ngay, nếu không chiếc xe không thể lưu hành. Cháu Hà liền nhận xét : “Sao lái xe không bảo giấy phạt đã bị gió thổi bay mất ?”. Anh Diệu mỉm cười, nhìn con nhận xét : “Con chỉ quen cái mẹo vặt dối trá mà thôi !.”. Còn tôi thầm nghĩ : “Cô bé này thông minh đấy”. Quả nhiên, sau này Hà đã trở thành TSKH, Phó GS toán học khi mới 30 tuổi.

GS. Phan Đình Diệu sinh năm 1936, thuộc thế hệ trí thức mới Việt Nam. Anh lại chuyên sâu vào toán học, nên không học chữ Hán, như người bạn rất thân của anh là GS. Sử học Hà Văn Tấn. Thế nhưng anh am hiểu niêm luật và sáng tác thể thơ Đường rất hay. Anh còn có một phẩm chất cao qúi của trí thức Việt Nam là không chỉ biết say mê nghiên cứu khoa học, mà còn quan tâm đến sự hưng thịnh của quốc gia. Vì thế anh hiểu biết, lo lắng đến sự phát triển của đất nước và đã đề xuất nhiều kiến giải có giá trị khoa học và thực tiễn. Sau ngày thống nhất 1975, trên các phương tiện thông tin nước ta thường sử dụng từ “  Hộp đen” để nói về sự chỉ đạo kinh tế. Anh và GS. Hoàng Tụy đã kịch liệt phê phán việc sử dụngtùy tiện thuật ngữ này. Tôi hỏi, tại sao nhà toán học lại quan tâm đến hai từ trên? Anh đã trả lời rằng, khái niệm “ Hộp đen” chỉ cho biết những thông tin đi vào, còn nó diễn biến ra sao bên trong thì không biết. Chỉ đạo nền kinh tế của một nước mà không biết được nó vận động như thế nào, đạt kết quả ra sao. . .thì thật vô cùng nguy hiểm, dễ đi đến thất bại, nên bọn mình phải có ý kiến phản bác! Trong những năm 1970, khi chế độ bao cấp còn rất triệt để, anh Diệu đề xuất là cần khai thác qui luật cạnh tranh, ví như trên cùng một trục đường nên có hai hợp tác xã mua bán, hai mậu dịch quốc doanh, thì chất lượng cung cấp hàng hoá sẽ được cải tiến, nhu cầu thị hiếu người mua được thoả mãn tốt hơn nhiều và giảm được hiện tượng bắt nẹt, cửa quyền đang khá phổ biến. Cũng trong những năm này, khi nhiều người còn hiểu lí thuyết thông tin anh đã kiến nghị nước nghiên cứu, xây dựng công nghệ thông tin. Nhớ lại một lần đến thăm ông Nguyễn Khắc Viện, tình cờ tôi gặp thư kí riêng của Tổng bí thư. Ông này cho biết, vị lãnh đạo cao nhất của Đảng có ý kiến gặp ông Viện và GS. Phan Đình Diệu. Sau đó ít lâu tôi tới nhà anh Diệu chơi, thời này đã dời về khu Đội Cấn, sực nhớ tôi bèn hỏi chuyện trên. Anh Diệu cho biết ông Tổng Bí thư có hỏi anh đánh giá như thế nào về trí thức Việt Nam. Anh trả lời rằng nước ta hiện chưa có trí thức. Tổng Bí thư ngạc nhiên hỏi lại: “ Còn giáo sư?”, anh trả lời, tôi cũng không phải trí thức. Tổng Bí thư ngạc nhiên thì anh cắt nghĩa rằng: trí thức là phải có chính kiến, còn bảo sao nghe vậy không phải là trí thức. Khi được hỏi một số nhà khoa học nổi tiếng của Việt Nam đang sinh sống và làm việc ở các nước phương Tây có phải là trí thức không, anh Diệu trả lời chưa phải, bởi lẽ họ sống xa Tổ quốc, nên chưa hiểu biết thực tế, để góp sức giải quyết khó khăn, trực tiếp đóng góp vào sự phát triển của đất nước.

Hôm về Hà Nội, xem triển lãm Giảng Võ, nơi gian hàng của Viện khoa học Việt Nam, tôi thấy anh trong bức ảnh Chủ tịch Hội đồng bộ trưởng Phạm Văn Đồng tiếp các nhà khoa học nước ngoài. Tiện gần nơi nhà anh ở, tôi ghé vào chơi. Cô Hương, vợ anh rót nước mời tôi, nói : “Anh chờ chốc lát, anh Diệu còn bận tay một tí”. Tôi thầm nghĩ chắc nhà khoa học đang mãi mê nghiên cứu, nên thông cảm vui vẻ ngồi chờ.  Khoảng gần phút sau anh Diệu đi ra, thân mật bắt tay tôi. Tôi hỏi: “ Bác đang say mê nghiên cứu gì vậy ?”. Anh hồn nhiên trả lời: “Mình đang làm cái chuồng gà”. Tôi ngạc nhiên thốt lên : “Chuồng gà bọn tôi làm mới đúng, còn bác phải dành thời giờ cho nghiên cứu, phát minh chứ!?”.  Anh Diệu trả lời : “Mình đang phải làm tiếp cái chuồng lợn nữa kia!”.

Vào khoảng đầu năm 1980, mấy học sinh lớp 9 trường cấp III Phan Đình Phùng niên khoá 1953-54 công tác ở Hà Nội tổ chức gặp mặt nhau. Tôi ở tỉnh lẻ, song tình cờ có việc về thủ đô, nên may mắn biết tin tới dự. Anh em đã đông đủ, chỉ chờ mỗi anh Diệu. Khoảng mười phút sau thấy anh tới, anh em đua nhau hỏi lí do. Anh trả lời : “Mình đi sớm, nhưng chiếc mobilet bị hỏng giữa đường”. Có một ai đó hỏi lại : “ Ô tô Viện trưởng đâu, sao không đi?”. Anh Diệu trả lời ngay : “Nhà nước cấp xe chỉ dùng vào việc công, việc riêng mình lấy dùng là sai!”. Trong cuộc họp mặt này có thầy học cũ của chúng tôi là GS. Bùi Văn Nguyên. Trước đấy, nhân đến tuổi thất thập, thầy làm bài thơ gửi mấy học trò cũ hồi học phổ thông hoạ. Tôi có làm một bài theo Đường luật, đưa anh Diệu xem và anh phát hiện có một câu sai luật. Anh Diệu là nhà toán học, song yêu thơ và làm thơ rất hay, vậy nên các bạn liền đề nghị anh đọc cho nghe một bài. Anh Diệu đã đọc bài thơ làm tặng Giáo sư, Viện sĩ Trần Đại Nghiã :

                                           “Nghĩa lớn gọi về với nước non

Buồn vui đã trải cuộc vuông tròn

Rèn tài văn võ thời phiêu bạt

Gánh việc giang sơn thuở mất còn

Tình nặng ấy chung tình đất nước

Nghiệp đời há kể chuyện vàng son

Gốc thông đứng thẳng dầu mưa nắng

Để gió lành reo mát nước non.”

Bài thơ này giới khoa học rất thích, nhiều báo chí đã đăng. Nghe nói khi GS. VS. Trần Đại Nghĩa sang Pháp công tác, có người  đọc cho nghe và ông đã khóc vì cảm động!

Một lần tôi cùng đi với GS. Hồ Ngọc Đại tới gặp anh Diệu để trao đổi về chương trình bộ môn toán các lớp bậc học phổ thông của Trung tâm thực nghiệm giáo dục, do GS. Đại làm Giám đốc. Anh Diệu đã không nể nang phê phán thẳng thừng nhiều vấn đề về chương trình môn toán thực nghiệm mà GS. Đại nêu ra. Do ý kiến bất đồng, nên không khí buổi gặp không được thoải mái. Mấy năm sau đó, một lần tôi nghe anh Diệu nhận xét: “ Hồ Ngọc Đại là người thực sự có năng lực”. GS. Đại tỏ ra rất vui vẻ khi nghe tôi thuật lại chuyện này.

Trong cương vị đại biểu Quốc hội, Uỷ viên thường vụ Mặt trận Tổ quốc Việt Nam, anh Diệu đã đề xuất nhiều ý kiến có cơ sở khoa học và đầy tâm huyết với Đảng và Nhà nước để xây dựng và phát triển đất nước. Đến những năm 1990, khi tin học, công nghệ thông tin ở nước ta bắt đầu phát triển, anh được chính phủ cử giữ cương vị Phó trưởng ban thường trực ban chỉ đạo chương trình quốc gia về công nghệ thông tin và là Chủ tịch đầu tiên của Hội tin học Việt Nam. Khi Chính phủ phê duyệt đề án 112 “ Tin học hoá hành chính Nhà nước”, nhằm mục đích xây dựng chính phủ điện tử từ giai đoạn 2001-2010, với tầm hiểu biết sâu về tin học và có một nhãn quan nhìn xa trông rộng, anh đã trình lên chính phủ một bản phân tích khoa học về nhiều mặt, để chứng minh nếu đề án 112 thực hiện sẽ không thể thành công. Tiếc rằng, vì những lí do nào đó, bản kiến nghị của anh đã không được chấp nhận. Quả nhiên, đến cuối năm 2005, đề án này đã thất bại vì có nhiều sai phạm và đến tháng 4 năm 2007, Thủ tướng chính phủ đã quyết định ngừng thực hiện.

Rất tiếc, từ tuổi 74 thì sức khoẻ anh không được tốt lắm. Song anh đã có hai niềm vui lớn. Đó là Công nghệ thông tin của việt Nam do anh và các cộng sự  đặt nền móng đã và đang phát triển nhanhchóng. Và, các con anh đã trưởng thành và đi theo con đường khoa học của anh.

Nhớ lại năm 1993, trường PTTH Phan Dình Phùng Hà Tĩnh tổ chức kĩ niệm 60 năm thành lập, là học sinh cũ, anh Diệu về dự. Trong lời  phát biểu, anh đã chúc học sinh của trường cố gắng học tập để có năng lực cống hiến thật nhiều cho quê hương, đất nước, như danh nho Nguyễn Công Trứ (1778-1858) từng nói :

Đã sinh ra ở trong trời đất

Phải có “danh” gì với núi sông.”

Cái “ danh” ở đây không phải là chức quyền, bổng lộc, tiền bạc. . ., mà là tạo ra nhiều sản phẩm tinh thần và vật chất cho đất nước, nhân dân. Trong suốt cuộc đời, Giáo sư TSKH. Phan Đình Diệu đã tâm niệm và làm như thế !